\documentclass[letterpaper,10.5pt]{article}

% --- PAQUETES BÁSICOS Y CONFIGURACIÓN REGIONAL ---
\usepackage[spanish,es-tabla]{babel}
\usepackage[margin=2.5cm,top=17mm,bottom=1.4cm,headheight=12pt]{geometry}
\usepackage{helvet}
\renewcommand{\familydefault}{\sfdefault}
\usepackage{microtype}
\usepackage{setspace}
\setstretch{1.2}
\setlength{\parindent}{0pt}
\setlength{\parskip}{5pt}

% --- COLORES CORPORATIVOS ---
\usepackage{xcolor}
\definecolor{primary}{HTML}{1A1D20}  
\definecolor{secondary}{HTML}{47607b}
\definecolor{divider}{HTML}{D0D5DD}
\definecolor{accent}{HTML}{005A9C}

% --- ENCABEZADOS Y PIES DE PÁGINA ---
\usepackage{fancyhdr}
\pagestyle{fancy}
\fancyhf{}
\renewcommand{\headrulewidth}{0.5pt}
\renewcommand{\headrule}{{\color{divider}\hrule width \headwidth height \headrulewidth}}
\renewcommand{\footrulewidth}{0.5pt}
\renewcommand{\footrule}{{\color{divider}\hrule width \headwidth height \headrulewidth}}

% Encabezado
\fancyhead[R]{%
    \raisebox{4mm}{%
        \parbox{\linewidth}{\raggedleft
            \fontsize{8.5pt}{13pt}\selectfont\bfseries\color{secondary}%
            CENTRO DE INVESTIGACIÓN, TECNOLOGÍA, EDUCACIÓN Y VINCULACIÓN ASTRONÓMICA (CITEVA)\\%
            \fontsize{8.5pt}{12pt}\selectfont\mdseries\color{secondary}%
            Universidad de Antofagasta --- Antofagasta, Chile
        }%
    }%
}

% Pie de página
\fancyfoot[C]{%
    \raisebox{-0.5mm}{%
    \fontsize{8.5pt}{12pt}\selectfont\color{secondary}%
    eduardo.unda@uantof.cl \ \textperiodcentered\ %
    +56 552 637 599 \ \textperiodcentered\ %
    astro.uantof.cl
    }%
}

% --- ELEMENTOS GRÁFICOS Y VÍNCULOS ---
\usepackage[colorlinks=true,linkcolor=accent,citecolor=accent,urlcolor=accent]{hyperref}

% --- INICIO DEL DOCUMENTO ---
\begin{document}
{\raggedleft\small\color{primary}%
\smash{\raisebox{-3mm}{Antofagasta, 21 de septiembre de 2026}}\par}

\vspace{0.5cm}

{\raggedright\small
\textbf{\color{primary} Señores} \\
\textbf{\color{primary} Comité Mixto ESO / Gobierno de Chile} \\
Convocatoria a Fondos Regionales 2026 \\
Presente\par}

\vspace{0.5cm}
{\centering\color{primary}\bfseries\fontsize{11.5pt}{15pt}\selectfont
CARTA DE PATROCINIO Y APOYO INSTITUCIONAL\par}
\vspace{0.3cm}

{\small\color{primary}
\textbf{Proyecto:} AstroGit-OS / Región de Antofagasta \\
\textbf{Director del Proyecto:} Pablo Zúñiga\par}

\vspace{0.3cm}
Estimados Miembros del Comité Mixto:
\vspace{0.3cm}

Por medio de la presente, manifiesto el interés y patrocinio institucional del \textbf{Centro de Investigación, Tecnología, Educación y Vinculación Astronómica (CITEVA)} de la Universidad de Antofagasta al proyecto \textbf{AstroGit-OS: Arquitectura Declarativa y Linaje Operacional para Instrumentación Astronómica mediante GitOps}, liderado por el ingeniero Pablo Zúñiga, en el marco de la Convocatoria 2026 de los Fondos Regionales del Comité Mixto ESO--Gobierno de Chile. Nuestra evaluación de la propuesta destaca su elevado valor científico, institucional y territorial para la Región de Antofagasta, destacando:

\begin{itemize}
    \vspace{-2.5mm}
    \item \textbf{Alineación estratégica en astroinformática de frontera} \\
    AstroGit-OS aborda un desafío crítico en mega-infraestructuras astronómicas: garantizar la trazabilidad, certidumbre y reproducibilidad operacional sin intervenir redes de control ni flujos de datos masivos, posicionando a la región como un polo de innovación en \textit{software} de vanguardia.
    \vspace{-1mm}
    \item \textbf{Formación de capital humano avanzado y sinergia académica} \\
    CITEVA valora la incorporación de talento local mediante la vacante de Asistente de Desarrollo en Astrofísica. Su articulación con programas de pre y posgrado universitario permitirá transferir estándares de ingeniería industrial y consolidar capacidades técnicas permanentes en la región.
    \vspace{-1mm}
    \item \textbf{Conocimiento abierto, transferencia tecnológica y proyección regional} \\
    El compromiso \textit{Open-Source} (licencias MIT y CC BY 4.0) y la realización de talleres técnicos en CITEVA fortalecen el ecosistema académico local, proyectando además la transferencia de esta arquitectura declarativa hacia la gestión de infraestructura en la gran minería regional.
\end{itemize}

Con base en estas fortalezas, CITEVA promoverá la difusión académica, la vinculación estudiantil y la coordinación técnica necesarias para el éxito del proyecto en la región. Consideramos que AstroGit-OS constituye un modelo ejemplar de desarrollo tecnológico regional con proyección científica e industrial, por lo que respaldamos su actual postulación a los Fondos Regionales ESO--Gobierno de Chile 2026.

\vspace{4mm}
Atentamente,
\vspace{18mm}

\vfill
\begin{center}
    \color {primary}
    \bfseries
    Dr. Eduardo Unda-Sanzana \\
    \mdseries
    Director \\
    Centro de Investigación, Tecnología, Educación y Vinculación Astronómica (CITEVA)\\
    Universidad de Antofagasta
\end{center}

\vspace{-2mm}

\end{document}